\indent

\textbf{Ethics} is the branch of philosophy that deals with what is right and wrong, good and bad, and how people should behave. It provides a framework for evaluating human actions, decisions, and intentions - especially when the right choice isn’t obvious.

In the professional IT context, Ethics refers to the \textbf{moral standards} that guide how IT professionals act when developing, managing, or using technology. It involves \textbf{responsible decision-making} - especially when those decisions affect people’s privacy, safety, or trust.

\subsection{Introduction}

\subsubsection{Key components of Ethics}

\indent 

\begin{itemize}
    \item Values: These are the deeply held beliefs or ideals that guide our decisions and shape what we consider right or wrong. Examples of core values include integrity, respect, responsibility, fairness, and trustworthiness.
    \item Principles: Principles are the general standards or rules derived from values that help us decide what actions are ethically acceptable. They provide structure - a moral compass that translates abstract values into concrete guidance. Examples of ethical principles in IT would be confidentiality, transparency, accountability, and non-maleficence.
    \item Behavior: Behavior is the observable expression of values and principles in real-life decisions and professional conduct. This is where ethics becomes visible - how professionals act when faced with a dilemma reveals their true ethical stance. Ethics is demonstrated through consistent behavior, not just belief.
    \item Consequences: Consequences are the results of our ethical (or unethical) actions. Evaluating consequences helps professionals understand whether their actions aligned with their intended ethical values.
\end{itemize}

\subsubsection{Our Ethical Framework}

\indent 

An \textbf{ethical framework} is a \textbf{structured way of thinking} about right and wrong. It helps professionals make \textbf{consistent, well-reasoned decisions}, especially when facing dilemmas where rules may be unclear or conflicting.

The IT world is full of \textbf{complex, high-impact decisions} - handling personal data, building AI models, managing automation, securing networks. Legal compliance alone isn’t enough; you often face situations where laws are silent or lag behind technology. A clear ethical framework ensures responsible innovation, public trust in technology, and professional accountability.

\subsubsection{Why Ethics matter in IT?}

\indent 

\begin{itemize}
    \item Technology shapes society: Information Technology now influences every part of modern life: communication, finance, healthcare, education, transportation, and even democracy. When IT professionals design systems, write algorithms, or process data, they’re shaping how people live and interact.
    \item IT Professionals handle sensitive information: Data is often described as the “new oil”, but it’s also deeply personal. Professionals in IT manage data that can reveal people’s habits, health, finances, and identities. Without ethical responsibility, that data can be misused, leaked, or sold without consent.
    \item Laws can't cover every situation: 
        \begin{itemize}
            \item Technology evolves faster than legislation. There are many situations where something is legal but not ethical, for example: Using publicly available data to profile users without consent. Automating layoffs using AI without human oversight.
            \item Ethics fills the gap between what is legal and what is right. It ensures professionals make responsible choices even when the law is silent or outdated.
        \end{itemize}
    \item Emerging technologies raise new ethical dilemmas
        \begin{itemize}
            \item As technologies like AI, Big Data, IoT, and automation advance, they raise novel ethical questions such as Who is responsible when an autonomous system fails? What biases are we embedding in AI systems?
            \item Ethics guides how we anticipate, design, and govern these technologies responsibly.
        \end{itemize}
\end{itemize}

\subsubsection{Laws vs Ethics}

\begin{table}[htbp]
\centering
\renewcommand{\arraystretch}{1.35}
\begin{tabular}{|p{3cm}|p{7cm}|p{7cm}|}
\hline
\textbf{Aspect} & \textbf{Laws} & \textbf{Ethics} \\
\hline

\textbf{Definition} &
\textbf{Laws are formal rules} established and enforced by governments or regulatory bodies. They define what is \textbf{legally permissible or punishable}. &
\textbf{Ethics refers to the moral principles and values} that govern an individual's or profession’s behavior. They go beyond legal compliance, asking what is \textbf{right, fair, and responsible}. \\
\hline

\textbf{Purpose} &
To maintain social order and justice &
To guide good and fair professional conduct \\
\hline

\textbf{Basis} &
Enforced by legal authority &
Driven by conscience and professional values \\
\hline

\textbf{Flexibility} &
Specific and fixed (can become outdated) &
Dynamic and adaptable to context \\
\hline

\textbf{Scope} &
Minimum acceptable standard of behavior &
Ideal standard of behavior (aspirational) \\
\hline

\end{tabular}
\caption{Comparison of Laws and Ethics}
\end{table}

\subsubsection{Core Ethical Principles in IT}

\indent 

Ethical principles act as moral guidelines that help IT professionals make responsible decisions when facing complex or ambiguous situations. They ensure that technology serves humanity fairly, safely, and transparently, aligning professional actions with both social good and individual responsibility.

The core principles are:

\begin{itemize}
    \item Integrity: Act honestly, consistently, and transparently in all professional dealings. In IT:
        \begin{itemize}
            \item Do not falsify data or manipulate analytics to satisfy management or clients.
            \item Accurately report project limitations, timelines, and risks.
            \item Maintain honesty even under pressure to deliver quick results.
        \end{itemize}
    \item Confidentiality and Privacy: Respect and protect the privacy and confidentiality of information entrusted to you. In IT:
        \begin{itemize}
            \item Protect sensitive data such as user credentials, financial records, or health information.
            \item Follow data minimisation principles - collect only what is necessary.
            \item Securely dispose of or anonymise data after use.
        \end{itemize}
    \item Fairness and Non-Discrimination: Treat all individuals and groups equitably, ensuring that technology does not perpetuate bias or inequality. In IT:
        \begin{itemize}
            \item Evaluate algorithms and datasets for bias or discrimination.
            \item Design accessible systems for all users, including those with disabilities.
            \item Avoid favoring certain users or clients unfairly.
        \end{itemize}
    \item Accountability and Responsibility: Accept responsibility for your actions, decisions, and the outcomes of the technology you develop or manage. In IT:
        \begin{itemize}
            \item Own up to mistakes and take corrective action quickly.
            \item Communicate potential risks or limitations to stakeholders.
            \item Follow due process for security disclosure or system failures.
        \end{itemize}
    \item Transparency and Honesty: Communicate clearly and openly about how systems work, what data they collect, and what limitations exist. In IT:
        \begin{itemize}
            \item Avoid hidden data practices or deceptive design (dark patterns).
            \item Inform users about how AI systems make decisions.
            \item Clearly document assumptions, models, and data sources.
        \end{itemize}
    \item Respect for Intellectual Property: Acknowledge and protect the ownership rights of others’ ideas, software, designs, and creative work. In IT:
        \begin{itemize}
            \item Do not copy or use others’ code, data, or algorithms without permission.
            \item Properly cite open-source material or licensed software.
            \item Respect software licenses and contribute ethically to collaborative projects.
        \end{itemize}
\end{itemize}

\subsubsection{ACS Code of Ethics}

\indent 

The Australian Computer Society (ACS) is the professional association 
for Australia’s ICT sector. All members are bound by the ACS Code of Professional Conduct, which establishes the ethical standards expected of ICT professionals in Australia.

The Code provides: A framework for ethical behavior; Guidance for decision-making when facing dilemmas; A basis for accountability to clients, employers, and society.

\begin{table}[htbp]
\centering
\renewcommand{\arraystretch}{1.35}
\begin{tabular}{|p{4.5cm}|p{10cm}|}
\hline
\multicolumn{2}{|c|}{\textbf{The 6 core values of the ACS Code of Ethics}} \\
\hline

\textbf{The Primacy of the Public Interest} &
Place the interests of the public above personal, business, or sectional interests. \\
\hline

\textbf{The Enhancement of Quality of Life} &
Strive to improve the quality of life of those affected by your work. \\
\hline

\textbf{Honesty} &
Be honest in your representation of skills, services, and products. \\
\hline

\textbf{Competence} &
Work competently and diligently for stakeholders. \\
\hline

\textbf{Professional Development} &
Enhance your own and others’ professional competence. \\
\hline

\textbf{Professionalism} &
Uphold and advance the honor and integrity of the profession. \\
\hline

\end{tabular}
\caption{ACS Code of Ethics: Six Core Values}
\end{table}

\subsection{Ethical Issues in IT}

\subsubsection{Privacy and Data Protection}

\indent 

Privacy and data protection involve ensuring that \textbf{personal, sensitive, or confidential data} is collected, stored, processed, and shared responsibly and lawfully.

Ethical principles involved: Primacy of public interest (ACS); Integrity, fairness, and accountability.

Ethical Questions:

\begin{itemize}
    \item Do users truly understand how their data is used?
    \item Should data be stored indefinitely “just in case”?
    \item Who is accountable for a data breach? the developer, manager, or company?
\end{itemize}

\subsubsection{Intellectual Property}

\indent

Intellectual property refers to \textbf{creations of the mind} - software, designs, algorithms, and other digital assets that are legally protected.

Ethical principles involved: Honesty, Competence, Professionalism

Ethical Questions:

\begin{itemize}
    \item Can you use AI-generated content that mimics another artist’s work?
    \item Should universities allow students to use ChatGPT for coding assignments?
\end{itemize}

\subsubsection{Cybersecurity responsibilities}

\indent

Cybersecurity ethics involve protecting \textbf{information systems and networks} against unauthorized access, misuse, or damage - while balancing privacy, freedom, and accountability.

Ethical principles involved: Public interest, Competence, Professionalism.

Ethical Questions:

\begin{itemize}
    \item Is it ethical to exploit a found vulnerability to “teach” a company a lesson?
    \item Do users have a right to know if their systems are unsafe?
\end{itemize}

\subsubsection{AI and Automation ethics}

\indent 

AI and automation ethics examine how intelligent systems \textbf{make decisions, impact society}, and \textbf{distribute responsibility} for outcomes.

Ethical principles involved: Public interest, Honesty, Quality of Life

Ethical Questions:

\begin{itemize}
    \item Who is accountable for AI’s actions? the developer or the algorithm?
    \item Should AI ever make life-and-death decisions (e.g., self-driving cars)?
\end{itemize}

\subsubsection{Workplace ethics in IT}

Workplace ethics covers professional behavior, communication, and conduct within IT organisations - ensuring respect, fairness, and accountability among colleagues and clients.

Ethical principles involved: Professionalism, Honesty, Enhancement of Quality of Life

Ethical Questions:

\begin{itemize}
    \item How do you respond when asked to “cut corners” on testing to meet deadlines?
    \item Should you report unethical behavior by a superior?
\end{itemize}

\subsection{Professional Code of Ethics}

\indent 

A \textbf{Professional Code of Ethics} is a \textbf{formal set of principles and standards} that guide the \textbf{behavior, decision-making, and conduct} of professionals within a specific field.

It defines what is considered \textbf{ethical and responsible behavior} and ensures that professionals act with \textbf{integrity, competence, and respect} toward clients, colleagues, and society.

The ICT industry impacts \textbf{privacy, safety, security, and equality} on a massive scale. A Professional Code of Ethics ensures that practitioners:

\begin{itemize}
    \item Use technology responsibly
    \item Avoid misuse of power or information
    \item Protect the public and environment
    \item Promote fairness, honesty, and respect
\end{itemize}

\subsubsection{ACS Code of Professional Conduct}

The \textbf{Australian Computer Society (ACS)} is Australia’s peak professional body for ICT practitioners. Its \textbf{Code of Professional Conduct} defines the \textbf{ethical and professional standards} expected from all ACS members - including students, graduates, and professionals.

The main purpose of the ACS Code is to ensure members:

\begin{itemize}
    \item Protect the \textbf{public interest} and promote social good
    \item Demonstrate \textbf{honesty, fairness, and integrity}
    \item Maintain \textbf{competence and accountability}
    \item Respect \textbf{privacy, confidentiality, and intellectual property}
    \item Promote \textbf{diversity, equity, and inclusion} in ICT workplaces
\end{itemize}

\begin{table}[htbp]
\centering
\renewcommand{\arraystretch}{1.35}
\begin{tabular}{|p{3cm}|p{7cm}|p{7cm}|}
\hline
\textbf{Aspect} & \textbf{ACS Code of Ethics} & \textbf{ACS Code of Professional Conduct} \\
\hline

\textbf{Nature} &
High-level moral principles &
Detailed professional rules \\
\hline

\textbf{Purpose} &
Guides ethical thinking and decision-making &
Guides enforceable professional behaviour \\
\hline

\textbf{Scope} &
Aspirational / Ideal &
Practical / Actionable \\
\hline

\textbf{Enforcement} &
Not strictly enforceable &
Enforceable; breaches can lead to sanctions \\
\hline

\textbf{Focus} &
Morality, public good, fairness &
Accountability, competence, professionalism \\
\hline

\textbf{Example} &
Anonymising user data to protect privacy &
Disclosing a security flaw per professional obligation \\
\hline

\end{tabular}
\caption{Comparison of ACS Code of Ethics and ACS Code of Professional Conduct}
\end{table}

\subsubsection{Personal vs Professional Ethics}

\indent 

Ethics can be applied at multiple levels.  In IT, it’s important to distinguish between personal ethics (your own moral compass) and professional ethics (standards required by your profession).

\begin{itemize}
    \item Personal Ethics: Personal ethics are an individual’s own moral values, principles, and beliefs about right and wrong. They are shaped by culture, religion, upbringing, and personal experiences.
    \item Professional Ethics: Professional ethics are the accepted standards of behavior for a specific profession. They are formalised in codes of conduct or codes of ethics (e.g., ACS Code of Professional Conduct).
\end{itemize}

Personal ethics influence professional ethics, but they must align with professional standards when on the job. Conflicts may arise, and professionals must balance personal beliefs with professional obligations.

Example:

Personal belief: “It’s fine to experiment on user data if it helps innovation.”

Professional obligation: ACS Code requires informed consent and privacy protection.

\subsection{Ethical Decision Making Models}

\indent 

\subsubsection{Consequentialism (Utilitarianism)}

Consequentialism is an ethical theory that \textbf{judges the morality of an action based on its outcomes or consequences}. The most well-known form of consequentialism is \textbf{Utilitarianism}.

Key features:

\begin{itemize}
    \item Outcome oriented: The \textbf{moral worth} of an action depends entirely on its results, not on intent or rules.
    \item Maximising good: Actions are judged by how much \textbf{overall good or utility} they produce (happiness, wellbeing, benefit).
    \item Impartiality: All affected parties’ interests are considered equally.
\end{itemize}

Strengths:

\begin{itemize}
    \item Practical: Focuses on real-world outcomes.
    \item Flexible: Can adapt to different contexts and technologies.
    \item Outcome-driven: Useful for risk-benefit analysis in IT projects.
\end{itemize}

Limitations:

\begin{itemize}
    \item Hard to predict all consequences accurately.
    \item May justify unethical means if the outcome seems beneficial (“ends justify the means”).
    \item Can conflict with personal or professional rules (e.g., privacy laws).
\end{itemize}

\subsubsection{Deontological (Duty-based)}

\indent 

Deontological ethics is an ethical theory that \textbf{focuses on the inherent rightness or wrongness of actions}, rather than their outcomes. It is also called \textbf{duty-based ethics}, because it emphasizes \textbf{adhering to moral rules, duties, or obligations}.

Key features:

\begin{itemize}
    \item Rule oriented: Actions are guided by \textbf{universal moral rules} or professional duties.
    \item Intent matters: The motive or intention behind the action is ethically significant.
    \item Consistency: Ethical rules apply \textbf{consistently to everyone}, independent of context or outcomes.
\end{itemize}

Strengths:

\begin{itemize}
    \item Provides \textbf{clear guidance}: rules are explicit and easy to follow.
    \item Emphasizes \textbf{moral integrity and accountability}.
    \item Aligns well with professional codes and legal compliance.
\end{itemize}

Limitations:

\begin{itemize}
    \item Can be rigid: may lead to outcomes that seem harmful or impractical.
    \item Conflicts may arise when duties clash (e.g., confidentiality vs. public safety).
    \item Less flexible for complex, real-world technological dilemmas.
\end{itemize}

\subsubsection{Virtue Ethics}

\indent 

Virtue Ethics is an ethical theory that \textbf{focuses on the character, moral virtues, and integrity of the individual}, rather than rules (deontology) or outcomes (consequentialism). It asks “\textbf{What kind of professional should I be}?” rather than just “What should I do?

Key features:

\begin{itemize}
    \item Character centric: Emphasizes cultivating \textbf{moral virtues} such as honesty, courage, integrity, fairness, and empathy.
    \item Habitual good conduct: Ethical behavior is a result of \textbf{developing good habits} over time.
    \item Practical wisdom: Professionals use judgment to navigate complex situations, balancing virtues appropriately.
\end{itemize}

Strengths:

\begin{itemize}
    \item Encourages \textbf{holistic moral development} of professionals.
    \item Flexible: considers context and nuance.
    \item Complements rules and outcomes - it promotes \textbf{ethical culture}, not just compliance.
\end{itemize}

Limitations

\begin{itemize}
    \item Less prescriptive: doesn’t provide explicit instructions for specific dilemmas.
    \item Can be \textbf{subjective} - what counts as a “virtue” may vary by culture or workplace.
    \item Requires \textbf{self-reflection and moral judgment}, which can be inconsistent.
\end{itemize}

\subsubsection{Ethics in Emerging technologies}

\indent 

Emerging technologies like \textbf{AI, Big Data, IoT, Cloud Computing, and Deepfakes} are reshaping society, but they bring \textbf{new ethical challenges}. ICT professionals must balance \textbf{innovation with responsibility}.

Ethical concerns:

\begin{itemize}
    \item Artificial Intelligence
        \begin{itemize}
            \item Algorithmic bias: AI may reinforce social inequalities.
            \item Transparency: “Black box” models make decisions difficult to explain.
            \item Accountability: Who is responsible if AI causes harm?
            \item Privacy: AI often requires large datasets that may include personal information.
        \end{itemize}
    \item Big Data
        \begin{itemize}
            \item Data privacy: Collection, storage, and sharing of massive datasets.
            \item Consent: Users may not understand how their data is used.
            \item Misuse: Predictive analytics could lead to discrimination (e.g., in hiring or insurance).
        \end{itemize}
    \item Internet of Things
        \begin{itemize}
            \item Security vulnerabilities: Devices can be hacked, exposing personal data.
            \item Surveillance: Continuous monitoring can infringe on privacy.
            \item Data ownership: Who controls the data collected by smart devices?
        \end{itemize}
    \item Cloud Computing
        \begin{itemize}
            \item Data sovereignty: Data stored across countries may violate local laws.
            \item Access control: Shared cloud resources may be exploited by malicious actors.
            \item Vendor responsibility: Cloud providers must uphold ethical and legal standards.
        \end{itemize}
    \item Deepfakes and Synthetic Media
        \begin{itemize}
            \item Misinformation: Deepfakes can be used to deceive or manipulate public opinion.
            \item Consent: Using someone’s likeness without permission is unethical.
            \item Trust erosion: Widespread deepfakes may undermine trust in digital content.
        \end{itemize}
\end{itemize}

\subsubsection{Balancing Innovation and Responsibility}

\indent 

Emerging technologies \textbf{offer immense benefits} but also \textbf{introduce risks}. ICT professionals must \textbf{apply ethical frameworks, ACS principles, and decision-making models} to: Protect users and society; Avoid harm and bias; Promote transparency and accountability.


Guiding questions:

\begin{itemize}
    \item Who benefits and who could be harmed?
    \item Are there fairness, privacy, or consent issues?
    \item Does this comply with professional codes and laws?
    \item Are the long-term societal consequences acceptable?
\end{itemize}

\subsection{Questions}

\subsubsection{End of Lecture Questions}

\begin{itemize}
    \item Define ethics and explain the difference between personal ethics and professional ethics.
    \item List and briefly describe the six key principles of the ACS Code of Professional Conduct.
    \item Identify five major ethical issues in IT and provide a real-world example for each.
    \item Explain a scenario where intellectual property rights could conflict with workplace demands.
    \item Compare and contrast Consequentialism, Deontological Ethics, and Virtue Ethics.
    \item When might personal ethics conflict with professional ethics? How should a professional resolve such a conflict?
\end{itemize}

\subsubsection{Case Study}

\indent 

You are a software engineer at a mid-sized Australian tech company developing an AI-powered recruitment platform. The system screens resumes and ranks candidates based on predicted job performance.

During testing, you discover that the AI algorithm systematically favors male candidates over female candidates, even though the data input was intended to be gender-neutral. Management insists that the platform be launched immediately because several large clients are waiting.

You are now faced with multiple considerations:

\begin{itemize}
    \item Public Interest: Launching a biased system could harm job applicants and damage public trust.
    \item Company Pressure: Delaying the launch may hurt client relationships and revenue.
    \item Professional Standards: The ACS Code of Professional Conduct requires fairness, honesty, and acting in the public interest.
    \item Legal Compliance: Discrimination in hiring is prohibited under Australian law.
\end{itemize}

\bigskip

\noindent\textbf{Case Study based Questions}

\begin{itemize}
    \item What are the main ethical issues in this scenario?
    \item Which ACS Code principles are most relevant to this situation?
    \item Using Consequentialism (Utilitarianism), what should you do? Consider both positive and negative consequences.
    \item Using Deontological Ethics, what action would be ethically required?
    \item Using Virtue Ethics, how should your personal and professional character guide your decision?
    \item Propose a step-by-step approach to address the bias in the AI system while balancing ethical, legal, and business concerns.
    \item How does this case illustrate the importance of professional accountability in IT?
\end{itemize}